% Backtest mechanics as implemented in code.
% Sources: backtesting/backtest_functions.py, lstm_models/lstm_common.py,
% itransformer_models/02_itransformer_backtesting.ipynb, lstm_models/02_lstm_backtesting.ipynb,
% lstm_models/03_lstm_diagrams.ipynb (metrics table).

\subsection{Backtest implementation}

Both model families are evaluated by the same code: the three strategies are imported
from \texttt{backtest\_functions.py} and the metrics, parameter search and evaluation
windows come from a single shared module, so the iTransformer and LSTM tables differ
only in the forecast that drives the signal.

\paragraph{Signal, entry and exit.}
The strategies act on a predicted percentage change. For the threshold and
volatility-adjusted strategies this is $\hat{r}_i = \hat{y}_i / P_{i-1}$, where
$\hat{y}_i$ is the model's forecast of the next hourly change; for the momentum
strategy it is
$\hat{r}_i = (P_{i-1} + \hat{y}_i - P_{i-L}) / P_{i-L}$ with lookback $L$, so the
forecast enters as a single additive term against a lookback move that is typically
orders of magnitude larger. A long is opened when $\hat{r}_i > \tau$, a short when
$\hat{r}_i < -\tau$, and no position otherwise. An open position is closed after
$H$ (\texttt{hold\_hours}) bars, or flipped immediately to the opposite side if the
reverse signal fires; a flip is recorded as an exit plus an entry and is charged the
transaction cost a second time. Positions are $\pm 1$ and fully symmetric; there is no
borrowing cost, margin model or liquidation logic.

\paragraph{Position size and leverage.}
The threshold strategy commits a constant $0.99$ of capital. The volatility-adjusted
and momentum strategies use
$w = \mathrm{clip}\!\left(\rho / \sigma_{i-1},\, 0.05,\, 1.00\right)$,
where $\rho$ is the tuned per-trade risk target and $\sigma$ is a 720-hour (30-day)
rolling standard deviation of hourly returns. The size is fixed at entry from the
volatility observed at $i-1$, so no look-ahead enters the sizing, and the cap of
$1.00$ means \textbf{leverage is never used}; minimum exposure is 5\%. Capital
compounds hourly as $C \leftarrow C\,(1 + s\,w\,r_i)$ with $s \in \{-1,0,1\}$ and
$r_i$ the realised hourly return.

\paragraph{Transaction costs.}
A single \texttt{transaction\_cost\_pct} of $0.0275\%$ is charged as $w \cdot C \cdot c$
on \textbf{entries only}. \textbf{No bid--ask spread, slippage,
funding rate, borrow fee or liquidation constraint is modelled anywhere in the code.}
For a symmetric long/short book on perpetual-style exposure this is a material
omission and the reported returns should be read as an upper bound.

\paragraph{Sharpe ratio and annualisation.}
Two different quantities appear, and they are not interchangeable.
\emph{(i)~The selection Sharpe} is annualised and computed on the hourly
portfolio-value path, including hours with no exposure:
$r_t = (V_t - V_{t-1})/V_{t-1}$, then
$\mathrm{SR} = \left(\bar{r} - r_f\right)/\mathrm{sd}(r)\times\sqrt{8760}$,
with the population standard deviation (\texttt{np.std}, $\mathrm{ddof}=0$) and
$8760$ hourly observations per year, which is correct for a market trading 24/7.
The \textbf{risk-free rate is $r_f = 0$} and is never overridden.
\emph{(ii)~The \emph{Sharpe} column of the backtesting results table} is a different
statistic: the unannualised mean-to-standard-deviation ratio of \emph{trade-to-trade}
portfolio returns taken from the trade log. It should not be compared with annualised
figures from the literature.

\paragraph{Parameter selection.}
Strategy parameters are tuned on the training period only, by a single Optuna study per
strategy and feature set whose objective is the \textbf{mean annualised Sharpe across ten
chronological folds} (\texttt{KFold}, no shuffling); the earlier procedure of running one
study per fold and averaging the per-fold winners was replaced, since averaging argmaxes
optimises nothing and is undefined for integer parameters. Search spaces:
$\tau\in[0.01,2.0]\%$, $H\in[1,168]\,\mathrm{h}$, $\rho\in[0,2.5]\%$,
$L\in[1,168]\,\mathrm{h}$; budgets of 200/500/1000 trials for the threshold,
volatility-adjusted and momentum strategies. Initial capital is \$10{,}000 in every run.
Because the study is executed with \texttt{n\_jobs=8}, Optuna's TPE sampler is not
deterministic even under a fixed seed, so repeated runs of identical code select
different parameters; strategy-level point estimates therefore carry search noise.

\paragraph{Regime windows.}
The uptrend and sideways periods are hard-coded constants
(2024-06-01--2024-12-01 and 2025-01-01--2025-07-01) in
\texttt{backtesting/data\_loader.py}. \textbf{No trend statistic, volatility filter or
regime-detection rule derives them from data}: they were chosen ex post as two
representative market phases. December 2024 falls in neither window, and the
forecasting stage uses slightly different windows than the backtest, so an ``uptrend''
sample in the forecasting table is not the same set of hours as in the backtesting
table.
